\section{Diagnóstico y solución de problemas}

\subsection{HomePilot no abre en el navegador}
\begin{lstlisting}[language=bash]
docker compose -f docker-compose.office.yml ps
curl http://127.0.0.1:3000/health
curl -I http://127.0.0.1:8080
\end{lstlisting}

Si la API se reinicia, consulte sus registros antes de reconstruir:
\begin{lstlisting}[language=bash]
docker compose -f docker-compose.office.yml logs --tail 150 homepilot-api
\end{lstlisting}

\subsection{Home Assistant no responde}
Primero verifique que el contenedor y el puerto de Home Assistant estén disponibles. En instalaciones existentes, HomePilot no debe reiniciar ni modificar Home Assistant sin autorización del responsable técnico.

\subsection{Error al enlazar el bridge}
Verifique la URL desde la red Docker, el token de larga duración y que el perfil seleccionado sea \texttt{bridge\_ha}. Si la URL funciona en el navegador, pero no desde HomePilot, puede ser un problema de red entre contenedores o una dirección \texttt{localhost} usada incorrectamente.

\subsection{Cámara demora o no muestra video}
Las cámaras pueden requerir varios segundos para establecer una sesión HLS o RTSP. Compruebe la disponibilidad en su integración de origen, la conectividad del equipo anfitrión y los registros de API. Evite abrir muchas transmisiones simultáneas en equipos con recursos limitados.

\subsection{Docker no puede descargar una imagen}
Un error como \texttt{502 Bad Gateway} desde Docker Hub suele ser transitorio o de conectividad. Verifique Internet, DNS, proxy corporativo y espacio disponible. Reintente la descarga; no borre volúmenes de datos como primera medida.

\subsection{Espacio insuficiente}
Use el diagnóstico de instalación para revisar imágenes, caché de compilación y volúmenes. La limpieza segura debe limitarse a caché y recursos no usados; preserve siempre las bases de datos y los volúmenes de servicios activos.

\begin{tcolorbox}[title=Escalación]
Antes de escalar un incidente a Nezu, adjunte: fecha y hora, perfil de instalación, salida de \texttt{--status}, estado de contenedores, últimos registros relevantes y una descripción exacta del síntoma. Nunca adjunte tokens ni contraseñas.
\end{tcolorbox}
